\chapter{解析结果}

\section{基态熵与约束层级}

在统计力学中，基态熵$S(0)=\ln\Omega_{0}$是基态简并度$\Omega_{0}$的自然对数（$k_{B}=1$），
反映了系统在零温下可访问的微观状态数。对于N皇后格点气体，基态精确对应于非攻击皇后构型，
因此$\Omega_{0}=Q(N)$。由Simkin的渐近公式，每皇后的基态熵为：
%
\begin{equation}
  s_{0} = \lim_{N\to\infty}\frac{\ln Q(N)}{N}
        = \ln N - \gamma,\quad \gamma = 1.94400(1).
  \label{eq:s0}
\end{equation}

这一表达式暗示了一个约束层级结构，如表\ref{tab:entropy}所示。
该层级提供了清晰的物理解释。从自由放置出发（每个皇后独立地从$N$列中选择一列，
总$\Omega=N^{N}$，每皇后熵为$\ln N$），首先施加列不可重复约束——此时每个皇后
必须选择不同的列，配置空间缩减为所有$N!$个列排列
（即``车问题''或rook problem\cite{Knuth2022}）。
由Stirling公式$\ln N!=N\ln N-N+\frac{1}{2}\ln(2\pi N)+O(1/N)$，每皇后熵变为：
%
\begin{equation}
  s_{\text{perm}} = \frac{\ln N!}{N}
  = \ln N - 1 + O\!\left(\frac{\ln N}{N}\right).
  \label{eq:sperm}
\end{equation}
%
在$N\to\infty$极限下，列约束恰好剪除1单位每皇后的熵——这就是Stirling修正；
有限$N$下还存在$O(\ln N/N)$的次导阶修正。物理上，这是$N$条等长列上
``恰好一个皇后''约束的联合效应。

\begin{table}[htbp]
  \centering
  \caption{逐级约束下每皇后的熵}
  \label{tab:entropy}
  \begin{tabular}{lcc}
    \toprule
    模型 & 约束 & 每皇后熵 $s=(1/N)\ln\Omega$ \\
    \midrule
    自由放置 & 仅行约束（每行一个皇后） & $\ln N$ \\
    排列（车问题） & 行$+$列（每列恰好一个） & $\ln N-1$ \\
    N皇后 & 行$+$列$+$两族对角线 & $\ln N-\gamma\approx\ln N-1.944$ \\
    \bottomrule
  \end{tabular}
\end{table}

进一步施加两族对角线约束后，每皇后熵从$\ln N-1$降至$\ln N-\gamma$，
额外代价为$\gamma-1\approx0.944$单位。大致而言，每族对角线约束对应的熵代价约为
$(\gamma-1)/2\approx0.472$单位——大约是列约束的一半。这一差异的根源在于
攻击线长度的不均匀分布：列有$N$条等长攻击线，对角线则有$2N-1$条不等长攻击线
（长度从1到$N$再到1）。短攻击线能承载的皇后数更少，因此其排斥约束相对更弱。

\section{高温极限}

当$T\to\infty$时，所有构型（无论是否有皇后互相攻击）变为等概率。
对于具有固定皇后数$N$的系统，构型总数为二项式系数$\binom{N^{2}}{N}$。
此时每皇后的熵简单地等于均匀分布的Shannon熵：
%
\begin{equation}
  \frac{S(\infty)}{N} = \frac{1}{N}\ln\binom{N^{2}}{N}.
  \label{eq:Sinf_def}
\end{equation}

利用Stirling近似$\ln n!=n\ln n-n+\frac{1}{2}\ln(2\pi n)+O(1/n)$，
当$N\gg1$时：
%
\begin{align}
  \ln\binom{N^{2}}{N}
  &= \ln(N^{2}!) - \ln N! - \ln(N^{2}-N)! \notag\\
  &= N\ln N - (N^{2}-N)\ln(1-1/N)
     + \Delta_{\mathrm{S}}.
  \label{eq:Sinf_stirling}
\end{align}
%
将$\ln(1-1/N)=-1/N-1/(2N^{2})-\cdots$展开，主导项给出
$-(N^{2}-N)\ln(1-1/N)=N-1/2+O(1/N)$，因此在大$N$极限下：
%
\begin{equation}
  \frac{S(\infty)}{N}
  = \ln N + 1 + O\!\left(\frac{\ln N}{N}\right).
  \label{eq:Sinf}
\end{equation}

高温熵包含两个主导项：自由放置贡献$\ln N$（与基态熵的领头项相同），
以及由棋盘有限尺寸带来的$+1$。基态熵$s_{0}=\ln N-\gamma$与高温熵
$S(\infty)/N=\ln N+1$之间的差值为$1+\gamma\approx2.944$。
这一差值恰好对应系统从$T=\infty$冷却至$T=0$过程中$C_{v}/T$的积分
——这正是热力学积分的核心内容。

对于高温下的平均能量，我们可以利用攻击线分解形式和期望的线性性精确求解。
对于包含$n_{\alpha}$个皇后的攻击线$\alpha$，其攻击对数为$\binom{n_{\alpha}}{2}$。
在$T\to\infty$、所有$\binom{N^{2}}{N}$个配置等概率的条件下，
可以计算每条攻击线上皇后数的分布。最终结果是，对于任意有限$N$，
高温极限下每皇后的精确平均能量为：
%
\begin{equation}
  \frac{\langle E\rangle}{N}\bigg|_{T\to\infty}
  = \frac{(5N-1)(N-1)}{3N(N+1)}.
  \label{eq:E_highT}
\end{equation}
%
当$N\to\infty$时，此表达式收敛到简洁的极限值$E/N\to5/3\approx1.667$。

值得注意的是，$5/3$这个结果可以从攻击线计数的几何论证中直观理解。
每条``攻击''由一对共享某条攻击线的皇后构成。对于均匀随机放置，
任何两个特定格点同时被占据的概率为$P=N(N-1)/[N^{2}(N^{2}-1)]\approx1/N^{2}$。
各攻击线上格点对的总数$S_{\text{tot}}$可以从行、列、两族对角线的几何关系中
精确求出：$S_{\text{tot}}=N(N-1)(5N-1)/3$。因此总攻击对数的期望为$P\cdot S_{\text{tot}}$，
除以$N$即得$E/N$。$N^{2}$和$N^{3}$项的精确抵消产生了简洁的常数结果。
